% 2021 全国甲卷理科 6：正视图与四个备选侧视图
% 依据固定上游 PDF 第 1 页。切口投影均落在方形边的中点；虚线/实线按遮挡关系区分。
\begin{tikzpicture}[font=\small,line width=0.4pt]
  \def\s{1.38}
  \newcommand{\cutleft}[1][]{\draw[#1] (0,{0.5*\s}) -- ({0.5*\s},0);}
  \newcommand{\cutright}[1][]{\draw[#1] ({0.5*\s},0) -- (\s,{0.5*\s});}
  % 正视图
  \begin{scope}[shift={({2.15*\s},{1.18*\s})}]
    \draw (0,0) rectangle (\s,\s);
    \cutleft[dashed]
    \node[anchor=south] at ({0.5*\s},\s+0.04) {正视图};
  \end{scope}
  % A：右下角、实线
  \begin{scope}[shift={(0,0)}]
    \draw (0,0) rectangle (\s,\s); \cutright
    \node[anchor=east] at (-0.05,{0.06*\s}) {A};
  \end{scope}
  % B：左下角、虚线
  \begin{scope}[shift={({1.55*\s},0)}]
    \draw (0,0) rectangle (\s,\s); \cutleft[dashed]
    \node[anchor=east] at (-0.05,{0.06*\s}) {B};
  \end{scope}
  % C：右下角、虚线
  \begin{scope}[shift={({3.10*\s},0)}]
    \draw (0,0) rectangle (\s,\s); \cutright[dashed]
    \node[anchor=east] at (-0.05,{0.06*\s}) {C};
  \end{scope}
  % D：左下角、实线
  \begin{scope}[shift={({4.65*\s},0)}]
    \draw (0,0) rectangle (\s,\s); \cutleft
    \node[anchor=east] at (-0.05,{0.06*\s}) {D};
  \end{scope}
\end{tikzpicture}
